% ============================================================
% INTRODUCTION — Full version for NeurIPS 2026
% ============================================================

\section{Introduction}
\label{sec:intro}

Discovering causal relationships from observational time series data is a
cornerstone problem in science. In climate science, understanding whether
El Ni\~no drives Indian Ocean Dipole variability informs seasonal
forecasting~\citep{runge2019pcmci}. In neuroscience, mapping directed
information flow between brain regions reveals the architecture of
cognition~\citep{seth2015granger}. In finance, identifying causal
dependencies between asset returns underpins systemic risk
analysis~\citep{granger1969investigating}.

The dominant framework for time series causal discovery---Granger
causality~\citep{granger1969investigating}---formalizes causation as
predictive utility: variable $X_j$ Granger-causes $X_i$ if the past of
$X_j$ helps predict $X_i$ beyond what $X_i$'s own past provides. While
conceptually clean, this framework faces three practical limitations.
First, the classical bivariate test assumes linearity and stationarity,
failing on the nonlinear, regime-switching dynamics common in real-world
systems. Second, multivariate extensions require $O(K^2)$ pairwise tests,
becoming prohibitively slow for large variable sets. Third, the test
provides a binary accept/reject decision with no natural way to rank
causal relationships by strength.

Modern extensions address some of these issues. PCMCI~\citep{runge2019pcmci}
combines the PC algorithm~\citep{spirtes2000causation} with momentary
conditional independence (MCI) testing to handle autocorrelation---a
pervasive challenge in time series. However, PCMCI inherits the
stationarity assumption, requires expensive conditional independence
tests, and scales as $O(K^2 \cdot \tau_{\max})$ with the number of lags
tested. Neural Granger causality~\citep{tank2021neural} replaces the
linear model with component-wise LSTMs (cLSTM) or MLPs (cMLP), using
group-lasso regularization on input weights to zero out non-causal
connections. This approach requires a specialized causal training
objective---distinct from standard prediction---and careful tuning of the
regularization parameter $\lambda$.

Concurrently, \emph{state space models} (SSMs) have emerged as a leading
architecture for sequence modeling. The Mamba
architecture~\citep{gu2023mamba} processes sequences through explicit
input-to-state ($\bB$) and state-to-output ($\bC$) projection matrices,
with the key innovation of making these projections data-dependent
(selective). This structured computation---distinct from the opaque
nonlinearities of transformers---raises a natural question: \emph{does a
Mamba model trained on multivariate time series prediction learn internal
representations that reflect the causal structure of its training data?}

We answer this question affirmatively and show that the connection is both
theoretically grounded and practically useful. Our key insight is
surprisingly simple: the product of the input and output weight matrices
of a trained Mamba model---$\bS = |\bW_{\text{out}} \cdot \bW_{\text{in}}|$---yields
a causal score matrix whose sparsity pattern recovers the Granger causal
graph. This extraction requires \emph{no causal-specific training}, no
regularization, and no statistical testing---only a single inspection of
the already-learned weights.

\paragraph{Why does this work?} A model trained to predict $\hat{X}_i(t+1)$
must learn which input variables $X_j(t)$ carry predictive information for
each output. The weight matrices encode this routing: $\bW_{\text{in}}$
determines which inputs the model attends to, and $\bW_{\text{out}}$
determines which outputs each hidden dimension contributes to. Their
product captures the net input-to-output information pathway---precisely
the quantity that Granger causality measures.

\paragraph{Contributions.} We make four contributions:
\begin{enumerate}
    \item \textbf{A novel connection} between SSM internals and causal
    structure. To our knowledge, no prior work has shown that trained SSM
    weight matrices encode Granger causal relationships
    (Section~\ref{sec:method}).

    \item \textbf{Theoretical guarantee}: For linear VAR($p$) processes,
    we prove that the SSM causal score matrix recovers the true Granger
    causal graph when the model has sufficient capacity
    (Theorem~\ref{thm:var_recovery}).

    \item \textbf{Comprehensive experiments}: On 6 synthetic and 4
    real-world benchmarks, \methodname{} achieves AUROC 0.70--1.00,
    outperforming PCMCI and cLSTM on all settings. On climate data,
    \methodname{} achieves AUROC 0.97, outperforming Granger (0.93);
    on financial data, it achieves perfect recovery
    (Section~\ref{sec:experiments}).

    \item \textbf{Training-free extraction}: Unlike cLSTM, our method works
    on \emph{any} pre-trained Mamba model without retraining, specialized
    losses, or regularization---enabling post-hoc causal analysis of
    existing models.
\end{enumerate}
